\section{Using and Extending \tool{}}
\label{sec:usage}

A typical session installs the two packages, builds a benchmark and runs one configuration:
\begin{center}
\begin{minipage}{0.95\textwidth}
\footnotesize
\begin{verbatim}
pip install -e . && pip install -e tsfm_lens
cd tsfm_benchmark && PYTHONPATH=. python3 example_runs/run_full.py \
       --config configs/example.yaml --out ./benchmark_out --references monash
cd ../tsfm_lens && python run.py --config configs/examples/medium_2model.yaml
\end{verbatim}
\end{minipage}
\end{center}
The CLI adds a preflight check (\texttt{-{}-doctor}), partial reruns (\texttt{-{}-stages}, \texttt{-{}-force}),
checks for new checkpoints or library upgrades (\texttt{-{}-check-alignment}, \texttt{-{}-check-adapter}) and an
environment diff (\texttt{-{}-verify-provenance}). A CPU-only mock configuration runs every stage in minutes in continuous
integration, and the worked example (\repo{tsfm_lens/docs/worked_example.md}) walks through a report. The seven-model analysis records
21.1 GPU-hours of stage time, mostly SAE training and the concept battery; stages through L3 take about 1.4 hours. Of eight
TSFM checkpoints probed, two (Timer, Time-MoE) load through the zero-code path and five with standalone packages
need a hand-written adapter (\fid{MN-25}), for which a template, a conformance check
(\texttt{-{}-check-adapter}) and a guide with a pitfall checklist (\repo{tsfm_lens/docs/ADDING_A_MODEL.md}) exist. New report sections declare their artifact requirements, new claims must use the
typed \texttt{Finding} record, and new configuration fields declare which artifacts they affect.

\section{Limitations}
\label{sec:limitations}

\paragraph{Models.} Seven transformer models with contiguous tokens; disjoint-lag tokens (Lag-Llama) get
reduced scope. Timer's features do not exceed chance, and with seven models, contrasts are confounded with architecture, size and
training data. The analysis is univariate: Chronos-2's cross-series group attention is not measured
(\fid{PM-13}), and Chronos-T5's decoder and Sundial's flow head are not captured.

\paragraph{Ground truth, external validity and confirmation.} Ground truth is exact only for the synthetic tier and
encodes our generators' assumptions; 12 of 33 concept parts are dominated by one real-derived generator, and pretraining
overlap cannot be verified. Real data are part of the analysis, not absent from it: 410 of the 965 dev series are real-derived
(Monash weather, ETT, and Chronos electricity and traffic) and the dictionaries are trained with 4{,}000 additional real
Monash windows. Raw real benchmarks are nonetheless a weaker bed for the causal questions: real rows give TimesFM
representations of lower effective dimension than the synthetic corpus (1.76 against 2.66) and more dead dictionary
atoms (\fid{MN-20}); weakly predictable real-derived series receive flat forecasts in 67--100\% of cases, where an
ablation has nothing to move (\fid{BM-06}); real series carry no generative ground truth; and public archives may be in
the models' pretraining data. A full run on a large raw benchmark such as GIFT-Eval remains to be done; the pipeline runs
on any sealed corpus. Held-out confirmation covers accuracy, geometry,
perturbation agreement, transfer, single-feature causal effects, atlas structure and failure prediction, on one private
split opened once (a second look needs a fresh epoch), and the 20 confirmed transfers come from two source concepts.

\paragraph{Effect sharing on the same series.} The cross-model causal claims are the ones that did not confirm, and the
known-answer test explains why: the shared-input rung is specific but detects a planted shared concept across
architectures in at most 0.4 of seeds, and more shared series do not help (\fid{MN-30}). We therefore built the redesign
this diagnosis suggested (ablating a concept's member subspace on each side, an equal-rank size-matched null, and
interval estimates against matched floors) and fixed its pass criterion in advance: sensitivity $\ge0.8$ on both planted
shared concepts, false agreement $\le0.05$, on held-out seeds. It also fails (sensitivity 0.5 and 0.375, with 0 of 14
false agreements), and the reason is a ceiling rather than a design detail: even the planted distributed subspace itself
does not move the forecast more than an equal-sized random removal (\fid{MN-32}). Per-series agreement of causal
effects across architectures is therefore not measurable by ablation at the effect sizes these models have, and
\tool{} compares effects across models at the level of concept profiles instead (Section~\ref{sec:rq2}).

\paragraph{SAE dependence, usability and completeness.} Concept results depend on dictionary size, sparsity,
revival and seed; the system measures this dependence but cannot remove it, and feature death tracks layer
geometry (alive atoms scale with effective dimension with log-log exponent 1.02, $\rho=0.58$ over 45 cells;
\fid{MP-05}). Readability for newcomers is a design intention; no user study measured comprehension. Controls catch only failures
for which a known-answer case exists.

\section{Conclusion}
\label{sec:conclusion}

TSFM interpretability is harder than it looks: tokens mean different things in different models, there is no
annotated ground truth, and causal effects are small enough that measurement artifacts dominate unless controlled.
\tool{} addresses this with a controlled benchmark, measured alignment, null- and reach-gated causal concepts and a
register-then-confirm protocol, behind one command and a report designed to be understandable by newcomers. On seven
TSFMs, causal features form a compact vocabulary centered on forecast level and seasonality; different architectures
select the same inputs for a concept, replicably across four sealed draws, while organizing their causal machinery by
model; causal importance sits in a minority of features that activation prominence does not reveal; and depth
signatures replicate held-out. Each of these conclusions is stated with its evidence class, and each held-out claim
was registered before the private data were opened.

\paragraph{Availability.} Code (MIT), configurations, the example run and the findings ledger are at
\url{\repourl}; the \texttt{dev} branch holds the full research log and study drivers.
